\section*{अध्याय \ref{ch:g11:prob} --- प्रायिकता और यादृच्छिक चर}

\begin{solution}{exo:g11:prob:1}
$52$ पत्तों पर \omterm{def:g11:prob:model}{एकसमान प्रतिरूप}। पान: $\frac{13}{52} = \frac14$। बादशाह: $\frac{4}{52} = \frac{1}{13}$। पान या बादशाह: $\frac{13 + 4 - 1}{52} = \frac{16}{52} = \frac{4}{13}$
(पान का बादशाह एक ही बार गिना जाता है)। न पान न बादशाह: $1 - \frac{4}{13} = \frac{9}{13}$।
\end{solution}

\begin{solution}{exo:g11:prob:2}
$S = 6$: युग्म $(1,5), (2,4), (3,3), (4,2), (5,1)$, इसलिए $\P(S = 6) = \frac{5}{36}$।

$S \geq 10$: $10$ के लिए तीन युग्म, $11$ के लिए दो, $12$ के लिए एक: $\P(S \geq 10) = \frac{6}{36} = \frac16$।

$S$ सम: योग $2, 4, 6, 8, 10, 12$ के $1 + 3 + 5 + 5 + 3 + 1 = 18$ युग्म हैं: \omterm{def:g10:proba:distribution}{प्रायिकता} $\frac{18}{36} = \frac12$।
\end{solution}

\begin{solution}{exo:g11:prob:3}
\omterm{def:g10:proba:distribution}{प्रायिकताओं} का योग $1$ है: $p = 1 - 0.3 - 0.2 - 0.4 = 0.1$। तब
\[
\E(X) = 0.3 \times (-1) + 0.2 \times 0 + 0.4 \times 2 + 0.1 \times 5
= -0.3 + 0.8 + 0.5 = 1 .
\]
\end{solution}

\begin{solution}{exo:g11:prob:4}
$\E(Y) = 2 \times 3 - 5 = 1$; $\V(Y) = 2^2 \times 4 = 16$; $\sigma(Y) = 4$।
\end{solution}

\begin{solution}{exo:g11:prob:5}
जीत $0$, $2$, $5$ यूरो हैं, जिनकी \omterm{def:g10:proba:distribution}{प्रायिकताएँ} $0.5$, $0.3$, $0.2$ हैं; $2$
यूरो का दाँव घटाने पर $G$ मान $-2, 0, 3$ लेता है:
\begin{center}
\begin{tabular}{c|ccc}
$g$ & $-2$ & $0$ & $3$\\
\hline
$\P(G=g)$ & $0.5$ & $0.3$ & $0.2$
\end{tabular}
\end{center}
$\E(G) = -1 + 0 + 0.6 = -0.4$: खिलाड़ी औसतन प्रति खेल $40$ सेंट गँवाता है --- प्रतिकूल।
\end{solution}

\begin{solution}{exo:g11:prob:6}
क्रम से निकालने पर: $\P(X = 2) = \frac35 \times \frac24 = \frac{3}{10}$ (लाल फिर लाल); $\P(X = 0) = \frac25 \times \frac14 = \frac{1}{10}$ (नीली फिर नीली); अतः $\P(X = 1) = 1 - \frac{3}{10} - \frac{1}{10} = \frac{6}{10}$। तब
\[
\E(X) = 0 \times \frac1{10} + 1 \times \frac6{10} + 2 \times \frac3{10}
= \frac{12}{10} = 1.2 ,
\]
जो अंतर्ज्ञान के अनुकूल है: दोनों निकालों में औसतन $\frac35$ लाल होती हैं।
\end{solution}

\begin{solution}{exo:g11:prob:7}
$\E(G^2) = 0.5 \times 4 + 0.3 \times 0 + 0.2 \times 9 = 3.8$, इसलिए संक्षिप्त सूत्र से $\V(G) = 3.8 - (-0.4)^2 = 3.8 - 0.16 = 3.64$ और $\sigma(G) \approx 1.9$ यूरो।
\end{solution}

\begin{solution}{exo:g11:prob:8}
लाभ $80$ में से दावा घटाकर मिलता है:
\begin{center}
\begin{tabular}{c|ccc}
$b$ & $80$ & $-920$ & $-4920$\\
\hline
$\P(B=b)$ & $0.94$ & $0.05$ & $0.01$
\end{tabular}
\end{center}
$\E(B) = 75.2 - 46 - 49.2 = -20$: इस दाम पर कंपनी औसतन प्रति बीमा-पत्र $20$ यूरो \emph{गँवाती} है।
ढके गए जोखिमों के लिए प्रीमियम बहुत कम है (धनात्मक प्रत्याशित लाभ के लिए उसे
$100$ यूरो से अधिक होना चाहिए)।
\end{solution}

\begin{solution}{exo:g11:prob:9}
यादृच्छिक अंक \omterm{def:g10:proba:distribution}{प्रायिकता} $\frac14$ से $+1$ और \omterm{def:g10:proba:distribution}{प्रायिकता} $\frac34$ से $x$ होते हैं:
\[
\E = \frac14 + \frac{3x}{4} = 0 \iff x = -\frac13 .
\]
दंड $-\frac13$ के साथ यादृच्छिक अनुमान औसतन कुछ नहीं देता।
\end{solution}

\begin{solution}{exo:g11:prob:10}
दो सिक्के: वे \omterm{def:g10:proba:distribution}{प्रायिकता} $\frac24 \times 2 =
\frac12$ से मेल खाते हैं (HH या TT)। खिलाड़ी का लाभ
बराबर \omterm{def:g10:proba:distribution}{प्रायिकताओं} से $+1$ या $-1$ है: दोनों खिलाड़ियों के लिए \omterm{def:g11:prob:expectation}{प्रत्याशा} $0$
--- खेल न्यायसंगत है।

तीन सिक्के: तीनों \omterm{def:g10:proba:distribution}{प्रायिकता} $\frac{2}{8} =
\frac14$ से मेल खाते हैं (HHH या TTT)। खिलाड़ी
\omterm{def:g10:proba:distribution}{प्रायिकता} $\frac14$ से $2$ पाता है और \omterm{def:g10:proba:distribution}{प्रायिकता} $\frac34$ से $1$ देता है: $\E = \frac24 - \frac34 = -\frac14$। खेल
प्रति दौर $25$ सेंट से ऐन के पक्ष में है; यदि वह तीनों के मेल पर $3$ यूरो
देती तो वह न्यायसंगत होता।
\end{solution}

\begin{solution}{exo:g11:prob:11}
\emph{1.} आयोजक $1000 \times 5 = 5000$ यूरो पाता है और कुल $2000 + 5 \times 200 + 20 \times 25 = 3500$ यूरो पुरस्कार में देता है:
लाभ \emph{अचर} $1500$ है (सारे टिकट बिक जाने के बाद कोई यादृच्छिकता नहीं
बचती)। किसी एक खिलाड़ी का लाभ $G$ (पुरस्कार घटा $5$-यूरो का टिकट):
\begin{center}
\begin{tabular}{c|cccc}
$g$ & $1995$ & $195$ & $20$ & $-5$\\
\hline\rule{0pt}{11pt}%
$\P(G=g)$ & $\dfrac{1}{1000}$ & $\dfrac{5}{1000}$ & $\dfrac{20}{1000}$
& $\dfrac{974}{1000}$
\end{tabular}
\end{center}

\emph{2.}
\[
\E(G) = \frac{1995 + 5 \times 195 + 20 \times 20 - 974 \times 5}{1000}
= \frac{1995 + 975 + 400 - 4870}{1000} = -1.5 .
\]
हर खिलाड़ी $1.50$ यूरो गँवाने की अपेक्षा करता है; $1000$ खिलाड़ियों पर यह $1500$
यूरो हुआ --- ठीक आयोजक का लाभ। बही-खाता बराबर है: खिलाड़ी औसतन जो गँवाते हैं
वही आयोजक बटोरता है।
\end{solution}

\begin{solution}{pb:g11:prob:1}
\textbf{1.} $G$ $-2, -1, 0, 1, 2, 3$ मान लेता है, और हर एक की \omterm{def:g10:proba:distribution}{प्रायिकता} $\frac16$ है: $\E(G) = \frac{-2 - 1 + 0 + 1 + 2 +
3}{6} = \frac12$
यूरो।

\textbf{2.} $\E(G^2) = \frac{4 + 1 + 0 + 1 + 4 + 9}{6} =
\frac{19}{6}$; $V(G) = \frac{19}{6} - \frac14 = \frac{35}{12}
\approx 2.92$; $\sigma(G) \approx 1.71$।

\textbf{3.} $\E(2G - 1) = 2 \times \frac12 - 1 = 0$; $V(2G - 1) = 4\,V(G) = \frac{35}{3}$।

\textbf{4.} $\E(\text{योग}) = 7$: न्यायसंगत दाम $7$ यूरो। $8$ पर प्रत्याशित हानि प्रति
खेल $1$ यूरो है।

\textbf{5.} न्यायसंगत प्रीमियम: $0.05 \times 800 = 40$ यूरो। $60$ लेने पर: प्रति बीमा-पत्र
प्रत्याशित लाभ $20$ यूरो।

\textbf{6.} $5:3$ के अनुपात में बाँटना \emph{बीते} को पुरस्कार देता है ---
पर दाँव उसका है जो \emph{आगे} जीते, और B अब भी जीत सकता है। सब कुछ A को दे
देना किसी संभावित जीत को निश्चित मान लेता है --- B की छोटी संभावना भी कुछ
मूल्य रखती है।

\textbf{7.} खेले जाने वाले दौर: $3$; परिणाम (हर दौर में A-जीत/B-जीत): AAA,
AAB, ABA, ABB, BAA, BAB, BBA, BBB। B शृंखला केवल तीनों जीतकर ले जाता है:
अकेला BBB। A $8$ में से $7$ स्थितियों में शृंखला जीतता है: न्यायसंगत
बँटवारा A के लिए $64 \times \frac78 = 56$ पिस्टोल और B के लिए $8$।

\textbf{8.} अतिरिक्त, निरर्थक दौर खेलने से किसी की जीत नहीं बदलती: A को उसकी
छठी जीत मिल जाने के बाद के दौर केवल औपचारिकता हैं। पर सीमा ठीक $3$ दौरों
पर तय कर देने से सभी $8$ परिणाम समसंभावी हो जाते हैं --- ईमानदार गिनती को
एक ही लंबाई की कहानियाँ चाहिए, और उन्हें भर देने में कुछ नहीं जाता।

\textbf{9.} $5$--$5$ पर: एक निर्णायक दौर, हर एक को $32$। $5$--$4$
पर: A को $\frac{64 + 32}{2} = 48$। $5$--$3$ पर: अगला दौर या तो $64$ तक ले जाता है (A उसे
जीत लेता है) या $5$--$4$ की स्थिति तक, जिसका मूल्य $48$ है: A का हिस्सा
$\frac{64 + 48}{2} = 56$ --- वही फ़र्मा का $56 : 8$, पीछे की ओर चलकर फिर निकाला हुआ।

\textbf{10.} A को $2$ चाहिए, B को $3$: $4$ काल्पनिक दौर खेलिए ($16$
परिणाम)। B शृंखला $3$ या $4$ B-जीतों के साथ ले जाता है: $4 + 1 = 5$ परिणाम।
बँटवारा: $B$: $64 \times \frac{5}{16} = 20$; $A$: $44$।

\textbf{11.} किसी अनिश्चित भावी प्राप्ति का न्यायसंगत वर्तमान मूल्य उसकी
\omterm{def:g11:prob:expectation}{प्रत्याशा} है --- जो आने वाला है उसका प्रायिकता-भारित औसत। बीमा (प्रश्न 5)
और वित्त (अनिश्चित प्राप्तियों का दाम लगाना) इसी वाक्य के उद्योग बन गए रूप
हैं; और मेज़ की दूसरी ओर से, जुआघर भी।

\textbf{12.} दाँव अंततः कहीं न कहीं जाता ही है: दोनों हिस्से ऐसे \omterm{def:g11:prob:rv}{यादृच्छिक
चर} हैं जिनका योग अचर $64$ है, इसलिए रैखिकता से $\E_A + \E_B = 64$ --- न्यायसंगत हिस्से
हर स्कोर पर दाँव को पूरा चुका देते हैं।

\textbf{13.} रैखिकता: $\E(X + Y) = 3.5 + 3.5 = 7$, बिना किसी सारणी के --- और यह जुड़े हुए पासों
के लिए भी सच है। एक पासे का \omterm{def:g11:prob:variance}{प्रसरण}: $\E(X^2) - 3.5^2 = \frac{91}{6} - \frac{49}{4} = \frac{35}{12}$; और स्वतंत्र पासों के लिए $V(X + Y) = \frac{35}{12} + \frac{35}{12} = \frac{35}{6}$।

\textbf{14.} हर नियत अतिथि अपनी टोपी \omterm{def:g10:proba:distribution}{प्रायिकता} $\frac1n$ से वापस पाता है (उसकी
टोपी $n$ में से एक है)। $n$ अतिथियों पर जोड़ने पर: $\E(X) = n \times \frac1n = 1$ --- औसतन एक
भाग्यशाली अतिथि, चाहे दावत में दो टोपियाँ हों या दो हज़ार।

\textbf{15.} $n = 2$: बँटवारे: दोनों सही ($X = 2$) या आपस में बदले हुए ($X = 0$):
$\E = 1$। $n = 3$: छहों बँटवारे $X = 3, 1, 1, 1, 0, 0$ देते हैं: $\E = \frac{6}{6} = 1$।

\textbf{16.} गणना ने $n$ अलग-अलग सूचकों (``अतिथि $i$ को अपनी टोपी
मिली'') की प्रत्याशाएँ जोड़ीं, और रैखिकता प्रत्याशाएँ \emph{बेशर्त} जोड़ती
है --- परावलंबन किसी गुणनफल या \omterm{def:g11:prob:variance}{प्रसरण} का सत्यानाश कर देता, पर प्रत्याशाओं
के योग का कभी नहीं। यही रैखिकता की महाशक्ति है: स्वतंत्रता के बारे में कोई
प्रश्न नहीं।

\textbf{17.} दोनों लाभ: $\E = 0.5 \times 10 - 5 = 0$ और $0.005 \times 1000 - 5 = 0$। (लाभ के) \omterm{def:g11:prob:variance}{प्रसरण}: A: $\E(G^2) = 0.5 \times 25 + 0.5 \times 25 = 25$; B: $0.005 \times 995^2 + 0.995 \times 25 \approx 4\,975$:
$\sigma_A = 5$, $\sigma_B \approx 70.5$। लोग B ख़रीदते हैं: वे \omterm{def:g11:prob:expectation}{प्रत्याशा} (जो दोनों में शून्य है) नहीं,
बल्कि \emph{प्रसार} ख़रीद रहे हैं --- जीवन बदल देने वाली दाहिनी पूँछ की एक
छोटी-सी संभावना।

\textbf{18.} उछाल $n$ पर पहली पट आने की \omterm{def:g10:proba:distribution}{प्रायिकता} $2^{-n}$ है और वह $2^n$
देती है: हर $n$ \omterm{def:g11:prob:expectation}{प्रत्याशा} में $2^{-n} \times 2^n = 1$ जोड़ता है, इसलिए \omterm{def:g11:prob:expectation}{प्रत्याशा} $1 + 1 + 1 + \dots$ है:
अनंत। फिर भी कोई समझदार व्यक्ति एक बार खेलने के लिए $1\,000$ नहीं देता:
\omterm{def:g10:proba:distribution}{प्रायिकता} $\frac{15}{16}$ से भुगतान अधिक से अधिक $16$ ही होता है। \omterm{def:g11:prob:expectation}{प्रत्याशा}
\emph{दोहराए जा सकने वाले} दाँवों के दीर्घ काल का सारांश है; जीवन में एक बार
खेला जाने वाला और जंगली पूँछ वाला खेल अपने \omterm{def:g11:stat:mean}{माध्य} से दाम नहीं पाता।

\textbf{19.} प्रत्याशाएँ: बीमा कराने पर $-20$; बिना बीमा $-0.01 \times 1000 = -10$। बिना बीमा
वाली रणनीति औसतन जीतती है --- उसके लिए जो हानि सह सके और दाँव बार-बार लगा
सके। जो परिवार $-1\,000$ के बाद टिक ही न सके वह औसत लेने के धंधे में है ही नहीं:
वह \omterm{def:g11:prob:variance}{प्रसरण} मिटाने के लिए $10$ अतिरिक्त देता है। बीमाकर्ता औसत लेते हैं;
परिवार निश्चितता ख़रीदते हैं।

\textbf{20.} जन्म दो अधीर जुआरियों के बीच दाँव बाँटते हुए; शक्ति रैखिकता से,
जो परावलंबन से अनुमति लिए बिना आशाएँ जोड़ लेती है; जोखिम के प्रति अंधी, जिसे
\omterm{def:g11:prob:variance}{प्रसरण} --- लॉटरियाँ, पीटर्सबर्ग का खेल, परिवार का बीमा-पत्र --- नापता है;
और अगले वर्ष राज्याभिषेक बृहत् संख्याओं के नियम से, जो बता देगा कि औसत सदा,
अंततः, वसूल क्यों कर लेते हैं।
\end{solution}
